% =====================================================================
% Sekcja 4 - Wyniki (wersja kompaktowa)
% Wartości liczbowe policzone bezpośrednio z plików w overnight_results/
% i results/ - są zsynchronizowane z notebookiem main_results_report.ipynb.
% =====================================================================

\section{Wyniki}
\label{sec:wyniki}

Wszystkie wyniki pochodzą z jednego przebiegu kodu z katalogu
\texttt{experiments/}, na ustalonych ziarnach losowości. Surowe
CSV z wynikami zapisane są w katalogu \texttt{overnight\_results/}
i można je odtworzyć z dołączonego notebooka
\texttt{main\_results\_report.ipynb}.

\subsection{Faza pilotażowa: Qwen-0,5B, czysty CoT, $n=150$}
\label{sec:wyniki-pilot}

Na 150 zadaniach (trzy ziarna po 50) model Qwen2.5-0,5B-Instruct
osiąga dokładność \textbf{44,7\%} - mieszczącą się w przedziale,
który dobrze nadaje się do badania korelacji pewności z poprawnością.

Już na tej skali widać pierwszy systematyczny wzorzec: średnia
entropia tokenowa jest lepszym predyktorem poprawności niż samo
perplexity ($r = -0{,}317$, $p < 0{,}001$ vs $r = -0{,}236$,
$p = 0{,}004$). Innymi słowy, sygnał z całego rozkładu
prawdopodobieństw niesie więcej informacji niż średnia liczona
tylko po tokenach, które model faktycznie wybrał. Ten sam wzorzec
powtarza się w każdym kolejnym ustawieniu.

\subsection{Agent ReAct z kalkulatorem: Qwen-1,5B, $n=60$}
\label{sec:wyniki-agent}

Na 60 zadaniach (dwa ziarna po 30) Qwen2.5-1,5B-Instruct
w pętli ReAct z pojedynczym narzędziem (kalkulator) osiąga
dokładność \textbf{28,3\%}. Spośród 60 trajektorii 12 (20\%)
zostaje obciętych - model albo przekracza limit pięciu wywołań
kalkulatora, albo łamie format dialogowy (np.\ zapomina o etykiecie
\texttt{Action:}).

Najciekawsza obserwacja w agencie dotyczy tego, \emph{gdzie} warto
mierzyć pewność. Metryki liczone tylko na segmentach \texttt{Action:} -
czyli tam, gdzie model formułuje wywołanie kalkulatora - dają
najsilniejszą korelację z poprawnością: \texttt{action\_ppl\_mean}
osiąga $r = -0{,}436$, $p = 0{,}002$, $n = 48$, z bootstrapowym
95\% CI $[-0{,}60,\ -0{,}24]$ (5000 resamplowań). Te same metryki
liczone tylko na finalnej odpowiedzi (\texttt{Answer:}) praktycznie
nie niosą żadnego sygnału ($r = -0{,}099$, $p = 0{,}50$, CI
$[-0{,}32,\ +0{,}21]$ - przedział ufności pokrywa zero).

Intuicja stojąca za tym wynikiem jest dość prosta. W momencie,
w którym model generuje finalną odpowiedź, wynik liczbowy znajduje
się już w jego kontekście - widoczny jako rezultat ostatniego
wywołania kalkulatora. Model w zasadzie kopiuje go z kontekstu,
więc jest bardzo pewny niezależnie od tego, czy całe rozumowanie
było prawidłowe. Pewność modelu w trakcie \emph{formułowania}
wywołań - gdy musi zdecydować, co policzyć - rozróżnia poprawne
i błędne odpowiedzi. Pewność na samej \emph{konkluzji} - już nie.

Sam rozkład pewności też różni się między typami segmentów.
Segmenty \texttt{Answer:} są skupione najmocniej (średnie PPL$=1{,}26$,
top-1 margin $=0{,}80$), segmenty \texttt{Action:} - pośrednio
(PPL$=1{,}46$, margin $=0{,}64$), segmenty \texttt{Thought:} - najmniej
(PPL$=1{,}59$, margin $=0{,}61$). Ten porządek utrzymuje się w obu
ziarnach.

\subsection{Kontrola czystego CoT: Qwen-1,5B, $n=60$}
\label{sec:wyniki-cot15}

Żeby zobaczyć, jak na te same pytania reaguje model bez pętli
agentowej, te same 60 zadań przepuszczono jeszcze raz przez
Qwen2.5-1,5B-Instruct, tym razem w czystym CoT. Dokładność rośnie
do \textbf{60,0\%} (wobec 28,3\% w wariancie agentowym), a najsilniejszy
sygnał korelacyjny w całym projekcie pojawia się właśnie tutaj:
średnia entropia tokenowa łańcucha rozumowania osiąga $r = -0{,}490$,
$p < 0{,}001$, $n = 60$. Samo perplexity też koreluje z poprawnością,
ale słabiej: $r = -0{,}283$, $p = 0{,}028$. Włożenie modelu w pętlę
ReAct nie likwiduje więc sygnału PPL/entropii - przenosi go z całej
kontynuacji na segmenty wywołań narzędzia.

\subsection{Porównanie zbiorcze}
\label{sec:headline}

Tabela~\ref{tab:headline} zestawia korelacje Pearsona między metrykami
a poprawnością w trzech głównych ustawieniach.

\begin{table}[h]
\centering
\small
\begin{tabular}{llrrr}
\toprule
\textbf{Ustawienie} & \textbf{Metryka} & $n$ & $r$ & $p$ \\
\midrule
0,5B CoT (pilot)        & \texttt{cot\_perplexity}           & 150 & $-0{,}236$ & $0{,}004$ \\
0,5B CoT (pilot)        & \texttt{cot\_mean\_token\_entropy} & 150 & $-0{,}317$ & $<0{,}001$ \\
1,5B CoT                & \texttt{cot\_perplexity}           & 60  & $-0{,}283$ & $0{,}028$ \\
1,5B CoT                & \texttt{cot\_mean\_token\_entropy} & 60  & $-0{,}490$ & $<0{,}001$ \\
1,5B Agent (cała traj.) & \texttt{agent\_ppl\_mean}          & 60  & $-0{,}354$ & $0{,}006$ \\
1,5B Agent (cała traj.) & \texttt{agent\_entropy\_mean}      & 60  & $-0{,}388$ & $0{,}002$ \\
1,5B Agent (action)     & \texttt{action\_ppl\_mean}         & 48  & $-0{,}436$ & $0{,}002$ \\
1,5B Agent (action)     & \texttt{action\_margin\_mean}      & 48  & $+0{,}454$ & $0{,}001$ \\
\bottomrule
\end{tabular}
\caption{Korelacje Pearsona między metrykami a poprawnością
w trzech ustawieniach ($n$ to liczba obserwacji po usunięciu
braków danych). Wartości policzone bezpośrednio z plików CSV
w katalogu \texttt{overnight\_results/}.}
\label{tab:headline}
\end{table}

Z tabeli wyłaniają się dwa powtarzające się wzorce, niezależne
od skali modelu i ustawienia. Po pierwsze, średnia entropia tokenowa
i top-1 margin konsekwentnie wypadają lepiej niż samo perplexity
jako predyktory poprawności - pełny rozkład prawdopodobieństw
niesie więcej informacji niż prawdopodobieństwa tylko wybranych
tokenów. Po drugie, wewnątrz agenta liczy się to, na której części
trajektorii mierzy się pewność: korelacje na segmentach
\texttt{action} są silniejsze niż na całych trajektoriach, a te
z kolei silniejsze niż na samym segmencie \texttt{answer}.
